\chapter{Wherein the Ideas of D'Artagnan  Begin To Clear Up a Little}

\lettrine{D}{'Artagnan} immediately took the offensive.  <Now that I have told you all, dear friend, or rather you have guessed  all, tell me what you are doing here, covered with dust and mud?>

\zz
Porthos wiped his brow, and looked around him with pride. <Why, it  appears,> said he, <that you may see what I am doing here.>

<No doubt, no doubt, you lift great stones.>

<Oh! to show these idle fellows what a man is,> said Porthos, with  contempt. <But you understand\longdash>

<Yes, that is not your place to lift stones, although there are many  whose place it is, who cannot lift them as you do. It was that which made me  ask you, just now. What are you doing here, baron?>

<I am studying topography, chevalier.>

<You are studying topography?>

<Yes; but you—what are you doing in that common dress?>

D'Artagnan perceived he had committed a fault in giving expression to his  astonishment. Porthos had taken advantage of it, to retort with a question.  <Why,> said he, <you know I am a bourgeois, in fact; my  dress, then, has nothing astonishing in it, since it conforms with my  condition.>

<Nonsense! you are a musketeer.>

<You are wrong, my friend; I have given in my resignation.>

<Bah!>

<Oh, mon Dieu! yes.>

<And you have abandoned the service?>

<I have quitted it.>

<You have abandoned the king?>

<Quite.>

Porthos raised his arms towards heaven, like a man who has heard extraordinary  news. <Well, that does confound me,> said he.

<It is nevertheless true.>

<And what led you to form such a resolution.>

<The king displeased me. Mazarin had disgusted me for a long time, as you  know; so I threw my cassock to the nettles.>

<But Mazarin is dead.>

<I know that well enough, parbleu! Only, at the period of his death, my  resignation had been given in and accepted two months. Then, feeling myself  free, I set off for Pierrefonds, to see my friend Porthos. I had heard talk of  the happy division you had made of your time, and I wished, for a fortnight, to  divide mine after your fashion.>

<My friend, you know that it is not for a fortnight my house is open to  you; it is for a year—for ten years—for life.>

<Thank you, Porthos.>

<Ah! but perhaps you want money—do you?> said Porthos, making  something like fifty louis chink in his pocket. <In that case, you  know\longdash>

<No, thank you; I am not in want of anything. I placed my savings with  Planchet, who pays me the interest of them.>

<Your savings?>

<Yes, to be sure,> said D'Artagnan: <why should I not  put by my savings, as well as another, Porthos?>

<Oh, there is no reason why; on the contrary, I always suspected  you—that is to say, Aramis always suspected you to have savings. For my  own part, d'ye see, I take no concern about the management of my  household; but I presume the savings of a musketeer must be small.>

<No doubt, relative to yourself, Porthos, who are a millionaire; but you  shall judge. I had laid by twenty-five thousand livres.>

<That's pretty well,> said Porthos, with an affable air.

<And,> continued D'Artagnan, <on the twenty-eighth of  last month I added to it two hundred thousand livres more.>

Porthos opened his large eyes, which eloquently demanded of the musketeer,  <Where the devil did you steal such a sum as that, my dear friend?>  <Two hundred thousand livres!> cried he, at length.

<Yes; which, with the twenty-five I had, and twenty thousand I have about  me, complete the sum of two hundred and forty-five thousand livres.>

<But tell me, whence comes this fortune?>

<I will tell you all about it presently, dear friend; but as you have, in  the first place, many things to tell me yourself, let us have my recital in its  proper order.>

<Bravo!> said Porthos; <then we are both rich. But what can I  have to relate to you?>

<You have to relate to me how Aramis came to be named\longdash>

<Ah! bishop of Vannes.>

<That's it,> said D'Artagnan, <bishop of Vannes.  Dear Aramis! do you know how he succeeded so well?>

<Yes, yes; without reckoning that he does not mean to stop there.>

<What! do you mean he will not be contented with violet stockings, and  that he wants a red hat?>

<Hush! that is promised him.>

<Bah! by the king?>

<By somebody more powerful than the king.>

<Ah! the devil! Porthos: what incredible things you tell me, my  friend!>

<Why incredible? Is there not always somebody in France more powerful  than the king?>

<Oh, yes; in the time of King Louis XIII it was Cardinal Richelieu; in  the time of the regency it was Cardinal Mazarin. In the time of Louis XIV it  is M\longdash>

<Go on.>

<It is M. Fouquet.>

<Jove! you have hit it the first time.>

<So, then, I suppose it is M. Fouquet who has promised Aramis the red  hat.>

Porthos assumed an air of reserve. <Dear friend,> said he,  <God preserve me from meddling with the affairs of others, above all from  revealing secrets it may be to their interest to keep. When you see Aramis, he  will tell you all he thinks he ought to tell you.>

<You are right, Porthos; and you are quite a padlock for safety. But, to  revert to yourself?>

<Yes,> said Porthos.

<You said just now you came hither to study topography?>

<I did so.>

<Tudieu! my friend, what fine things you will do!>

<How do you mean?>

<Why, these fortifications are admirable.>

<Is that your opinion?>

<Decidedly it is. In truth, to anything but a regular siege, Belle-Isle  is absolutely impregnable.>

Porthos rubbed his hands. <That is my opinion,> said he.

<But who the devil has fortified this paltry little place in this  manner?>

Porthos drew himself up proudly: <Did I not tell you who?>

<No.>

<Do you not suspect?>

<No; all I can say is that he is a man who has studied all the systems,  and who appears to me to have stopped at the best.>

<Hush!> said Porthos; <consider my modesty, my dear  D'Artagnan.>

<In truth,> replied the musketeer, <can it be  you—who—oh!>

<Pray—my dear friend\longdash>

<You who have imagined, traced, and combined between these bastions,  these redans, these curtains, these half-moons; and are preparing that covered  way?>

<I beg you\longdash>

<You who have built that lunette with its retiring angles and its salient  edges?>

<My friend\longdash>

<You who have given that inclination to the openings of your embrasures,  by means of which you so effectively protect the men who serve the guns?>

<Eh! mon Dieu! yes.>

<Oh! Porthos, Porthos! I must bow down before you—I must admire  you! But you have always concealed from us this superb, this incomparable  genius. I hope, my dear friend, you will show me all this in detail.>

<Nothing more easy. Here lies my original sketch, my plan.>

<Show it me.> Porthos led D'Artagnan towards the stone that  served him for a table, and upon which the plan was spread. At the foot of the  plan was written, in the formidable writing of Porthos, writing of which we  have already had occasion to speak:—

<Instead of making use of the square or rectangle, as has been done to  this time, you will suppose your place inclosed in a regular hexagon, this  polygon having the advantage of offering more angles than the quadrilateral  one. Every side of your hexagon, of which you will determine the length in  proportion to the dimensions taken upon the place, will be divided into two  parts, and upon the middle point you will elevate a perpendicular towards the  centre of the polygon, which will equal in length the sixth part of the side.  By the extremities of each side of the polygon, you will trace two diagonals,  which will cut the perpendicular. These will form the precise lines of your  defense.>

<The devil!> said D'Artagnan, stopping at this point of the  demonstration; <why, this is a complete system, Porthos.>

\begin{bwbigpic}
	[\basicscale] 
	{70system}
	[This is a complete system]
\end{bwbigpic}

<Entirely,> said Porthos. <Continue.>

<No; I have read enough of it; but, since it is you, my dear Porthos, who  direct the works, what need have you of setting down your system so formally in  writing?>

<Oh! my dear friend, death!>

<How! death?>

<Why, we are all mortal, are we not?>

<That is true,> said D'Artagnan; <you have a reply for  everything, my friend.> And he replaced the plan upon the stone.

But however short the time he had the plan in his hands, D'Artagnan had  been able to distinguish, under the enormous writing of Porthos, a much more  delicate hand, which reminded him of certain letters to Marie Michon, with  which he had been acquainted in his youth. Only the India-rubber had passed and  repassed so often over this writing that it might have escaped a less practiced  eye than that of our musketeer.

<Bravo! my friend, bravo!> said D'Artagnan.

<And now you know all that you want to know, do you not?> said  Porthos, wheeling about.

<Mordioux! yes, only do me one last favour, dear friend!>

<Speak, I am master here.>

<Do me the pleasure to tell me the name of that gentleman who is walking  yonder.>

<Where, there?>

<Behind the soldiers.>

<Followed by a lackey?>

<Exactly.>

<In company with a mean sort of fellow, dressed in black?>

<Yes, I mean him.>

<That is M. Getard.>

<And who is Getard, my friend?>

<He is the architect of the house.>

<Of what house?>

<Of M. Fouquet's house.>

<Ah! ah!> cried D'Artagnan, <you are of the household  of M. Fouquet, then, Porthos?>

<I! what do you mean by that?> said the topographer, blushing to  the top of his ears.

<Why, you say the house, when speaking of Belle-Isle, as if you were  speaking of the château of Pierrefonds.>

Porthos bit his lip. <Belle-Isle, my friend,> said he,  <belongs to M. Fouquet, does it not?>

<Yes, I believe so.>

<As Pierrefonds belongs to me?>

<I told you I believed so; there are no two words to that.>

<Did you ever see a man there who is accustomed to walk about with a  ruler in his hand?>

<No; but I might have seen him there, if he really walked there.>

<Well, that gentleman is M. Boulingrin.>

<Who is M. Boulingrin?>

<Now we are coming to it. If, when this gentleman is walking with a ruler  in his hand, any one should ask me,—<who is M. Boulingrin?> I  should reply: <He is the architect of the house.> Well! M. Getard  is the Boulingrin of M. Fouquet. But he has nothing to do with the  fortifications, which are my department alone; do you understand? mine,  absolutely mine.>

<Ah! Porthos,> cried D'Artagnan, letting his arms fall as a  conquered man gives up his sword; <ah! my friend, you are not only a  Herculean topographer, you are, still further, a dialectician of the first  water.>

<Is it not powerfully reasoned?> said Porthos: and he puffed and  blew like the conger which D'Artagnan had let slip from his hand.

<And now,> said D'Artagnan, <that shabby-looking man,  who accompanies M. Getard, is he also of the household of M. Fouquet?>

<Oh! yes,> said Porthos, with contempt; <it is one M.  Jupenet, or Juponet, a sort of poet.>

<Who is come to establish himself here?>

<I believe so.>

<I thought M. Fouquet had poets enough, yonder—Scudery, Loret,  Pelisson, La Fontaine? If I must tell you the truth, Porthos, that poet  disgraces you.>

<Eh!—my friend; but what saves us is that he is not here as a  poet.>

<As what, then, is he?>

<As printer. And you make me remember, I have a word to say to the  cuistre.>

<Say it, then.>

Porthos made a sign to Jupenet, who perfectly recollected D'Artagnan, and  did not care to come nearer; which naturally produced another sign from  Porthos. This was so imperative, he was obliged to obey. As he approached,  <Come hither!> said Porthos. <You only landed yesterday and  you have begun your tricks already.>

<How so, monsieur le baron?> asked Jupenet, trembling.

<Your press was groaning all night, monsieur,> said Porthos,  <and you prevented my sleeping, corne de bœuf!>

<Monsieur\longdash> objected Jupenet, timidly.

<You have nothing yet to print: therefore you have no occasion to set  your press going. What did you print last night?>

<Monsieur, a light poem of my own composition.>

<Light! no, no, monsieur; the press groaned pitifully beneath it. Let it  not happen again. Do you understand?>

<Yes, monsieur.>

<You promise me?>

<I do, monsieur!>

<Very well; this time I pardon you. Adieu!>

The poet retreated as humbly as he had approached.

<Well, now we have combed that fellow's head, let us  breakfast.>

<Yes,> replied D'Artagnan, <let us breakfast.>

<Only,> said Porthos, <I beg you to observe, my friend, that  we only have two hours for our repast.>

<What would you have? We will try to make two hours suffice. But why have  you only two hours?>

<Because it is high tide at one o'clock, and, with the tide, I am  going to Vannes. But, as I shall return to-morrow, my dear friend, you can stay  here; you shall be master; I have a good cook and a good cellar.>

<No,> interrupted D'Artagnan, <better than that.>

<What?>

<You are going to Vannes, you say?>

<To a certainty.>

<To see Aramis?>

<Yes.>

<Well! I came from Paris on purpose to see Aramis.>

<That's true.>

<I will go with you then.>

<Do; that's the thing.>

<Only, I ought to have seen Aramis first, and you after. But man  proposes, and God disposes. I have begun with you, and will finish with  Aramis.>

<Very well!>

<And in how many hours can you go from here to Vannes?>

<Oh! pardieu! in six hours. Three hours by sea to Sarzeau, three hours by  road from Sarzeau to Vannes.>

<How convenient that is! Being so near to the bishopric; do you often go  to Vannes?>

<Yes; once a week. But, stop till I get my plan.>

Porthos picked up his plan, folded it carefully, and engulfed it in his large  pocket.

<Good!> said D'Artagnan aside; <I think I now know the  real engineer who is fortifying Belle-Isle.>

Two hours after, at high tide, Porthos and D'Artagnan set out for  Sarzeau.  